\chapter{L'empirisme de l'impact sur les prix}\label{ch:mx:the-empirics-of-price-impact}

Acheter un pour cent du volume d'une journée fait bouger le cours, en moyenne, d'un certain montant~; acheter quatre pour cent le fait bouger à peu près deux fois plus, et non quatre fois plus.
Que l'impact d'un ordre croisse comme la racine carrée de sa taille est l'une des régularités les plus robustes de la finance, et personne ne l'a conçue. Elle est aussi
difficile à mesurer~: l'impact est une petite moyenne noyée sous un grand bruit, les ordres mesurés ont été envoyés pour des raisons qui font bouger les cours eux aussi, et un
marché simulé ne peut montrer d'impact que s'il réagit. Ce chapitre mesure l'impact à trois échelles, une transaction isolée, le flux d'un intervalle et un métaordre, dans
un marché simulé réactif, puis ajuste la loi en racine carrée sur les données de métaordres de la firme et montre comment une prévision de leur calendrier la fausse.

\section{L'impact d'une transaction isolée}

L'impact est le mouvement de cours qui suit une transaction dans son sens (livre~7, chapitre~18)~; le chapitre~5 l'a mesuré par transaction, par rapport au milieu, à un
horizon donné. Le marché du chapitre est la population \texttt{firm.agentmkt} sur \texttt{firm.exchsim}~: des fournisseurs de liquidité qui placent des ordres d'un lot par rapport au
carnet tel qu'il est, une part à un tick du cours vendeur opposé et le reste jusqu'à dix ticks plus loin, avec un poids croissant à mesure que l'on s'éloigne (un carnet dont la
profondeur croît à mesure qu'on s'éloigne du prix), des preneurs de bruit dont les ordres arrivent à des instants tirés à l'avance, avec des signes et des tailles tirés à l'avance, et des fondamentalistes qui échangent vers une valeur cachée
qui se déplace par sauts aléatoires. Tout sauf les ordres de bruit et les sauts de la valeur réagit au carnet.

\begin{definition}[Courbe d'impact]\label{def:mx:the-empirics-of-price-impact:curve}
Une \emph{courbe d'impact}\index{courbe d'impact} est le mouvement moyen du cours dans le sens d'une transaction, ou d'un métaordre, en fonction du temps écoulé depuis
son début (la fonction de réponse d'une transaction isolée, en transactions ou en secondes) ou de sa taille.
\end{definition}

La réponse à un \omterm{def:mx:the-limit-order-book:orders}{ordre au marché} isolé, agrégé à partir de ses exécutions, croît avec sa taille sur l'heure simulée (763\,900 actions en 7\,639 transactions)~: après un
ordre supplémentaire, elle vaut 0.27 tick pour un lot, 0.58 pour deux ou trois, 1.12 pour quatre à sept et 2.34 pour huit ou plus (\cref{fig:mx:the-empirics-of-price-impact:single},
à gauche). Elle décroît ensuite pour les gros ordres, jusqu'à 0.46 tick cinquante ordres plus tard, tandis qu'elle croît pour les ordres d'un lot, jusqu'à 0.81~: les fondamentalistes qui connaissent la
valeur envoient des ordres d'un lot, et leurs transactions prédisent le cours. Dans ce carnet mince, la réponse est presque proportionnelle à la taille~; sur les marchés réels, elle est
fortement concave, en partie parce que les gros \omterm{def:mx:the-limit-order-book:orders}{ordres au marché} sont envoyés lorsque le carnet peut les absorber.

\begin{definition}[Impact agrégé]\label{def:mx:the-empirics-of-price-impact:aggregate}
L'\emph{impact agrégé}\index{impact agrégé} est la variation moyenne du cours sur un intervalle, en fonction du flux d'ordres signé de cet intervalle (achats moins
ventes), quels que soient les acteurs qui ont échangé.
\end{definition}

Sur des intervalles de dix secondes, le milieu bouge de 2.1 ticks par millier d'actions d'achat net~; les moyennes par classe vont de $-3.1$ ticks pour environ $-1\,400$ actions à
$+3.1$ pour $+1\,300$ (\cref{fig:mx:the-empirics-of-price-impact:single}, à droite). L'\omterm{def:mx:the-empirics-of-price-impact:aggregate}{impact agrégé} est linéaire pour de petits déséquilibres et s'infléchit pour les grands,
là où le flux rencontre les niveaux plus profonds du carnet et les fondamentalistes.

\begin{omfigure}
\begin{tikzpicture}
\begin{groupplot}[group style={group size=2 by 1,horizontal sep=1.6cm},width=6.6cm,height=5cm,
  tick label style={font=\footnotesize},label style={font=\small},grid=major,grid style={gray!25},table/col sep=comma]
\nextgroupplot[title={\small réponse à un ordre (ticks)},xmode=log,xlabel={ordres plus tard},ymin=0,ymax=2.6,
  legend columns=2,legend style={font=\footnotesize,at={(0.5,-0.3)},anchor=north,draw=none,fill=none}]
\addplot[omDef,very thick,mark=*] table[x=lag,y=s0]{figdata/microstructure/11-the-empirics-of-price-impact/response.csv};
\addplot[omThm,very thick,mark=square*] table[x=lag,y=s1]{figdata/microstructure/11-the-empirics-of-price-impact/response.csv};
\addplot[omProp,very thick,mark=triangle*] table[x=lag,y=s2]{figdata/microstructure/11-the-empirics-of-price-impact/response.csv};
\addplot[black,very thick,mark=diamond*] table[x=lag,y=s3]{figdata/microstructure/11-the-empirics-of-price-impact/response.csv};
\legend{1 lot,{2 à 3 lots},{4 à 7 lots},{8 lots et plus}}
\nextgroupplot[title={\small variation du milieu sur 10 secondes (ticks)},xlabel={volume signé (actions)},xtick={-2000,0,2000},
  scaled x ticks=false,xticklabel style={/pgf/number format/.cd,fixed,1000 sep={\,}}]
\addplot[omDef,very thick,mark=*] table[x=flow,y=dm]{figdata/microstructure/11-the-empirics-of-price-impact/aggregate.csv};
\end{groupplot}
\end{tikzpicture}

\omcaption{Impact à deux échelles dans le marché simulé. À gauche~: la variation moyenne du milieu dans le sens d'un \omterm{def:mx:the-limit-order-book:orders}{ordre au marché}, selon sa taille, en fonction du nombre d'ordres
écoulés depuis. À droite~: la variation du milieu sur des intervalles de dix secondes en fonction de leur volume signé net, regroupé par classes. Données~: \texttt{mx\_impact.single\_and\_aggregate}.}\label{fig:mx:the-empirics-of-price-impact:single}
\end{omfigure}

\section{L'impact d'un métaordre~: pendant et après}

\begin{definition}[Impact maximal et retour de l'impact]\label{def:mx:the-empirics-of-price-impact:peak}
L'\emph{impact maximal}\index{impact maximal} d'un métaordre est le mouvement du cours dans son sens, depuis son début jusqu'à sa dernière exécution. Le \emph{retour de l'impact}\index{retour de l'impact}
(reversion) est la part de ce mouvement qui disparaît après la fin du métaordre~; ce qui subsiste après un long moment est son
impact permanent.
\end{definition}

L'impact d'un métaordre est de quelques ticks sous un cours qui oscille de plusieurs ticks par minute, de sorte qu'une seule exécution en dit peu. Un simulateur peut faire mieux
que les données~: on fait tourner la séance deux fois, avec et sans le métaordre, avec le même bruit d'ordres et les mêmes sauts de valeur tirés à l'avance, puis on prend la différence des
deux trajectoires de cours. Le bruit qui ne dépend pas du métaordre s'annule~; ce que les fournisseurs et les fondamentalistes font différemment ne s'annule pas, car une petite
modification du carnet change leurs choix ultérieurs.

Deux résultats en découlent. D'abord, un fond rejoué ne peut montrer aucun impact. Injecté dans le flux rejoué de \texttt{firm.tape}, un métaordre de 24\,000 actions
laisse la trajectoire du milieu exactement où elle était sans lui, à tous les horizons de 5 à 20 minutes~: les fournisseurs de liquidité rejoués continuent d'afficher aux prix
qu'ils avaient affichés dans l'enregistrement et n'apprennent jamais l'existence de l'ordre. Ensuite, dans le marché réactif (\cref{fig:mx:the-empirics-of-price-impact:meta}), 30 fragments
sur dix minutes d'un achat de 24\,000 actions, soit 3.2~\% des 749\,200 actions de l'heure, font monter le cours de 6.3 ticks à mi-parcours (erreur type 1.9 sur six séances)
et de 7.8 (2.2) à la fin~; cinq minutes après la fin, il vaut $-1.5$ (1.5)~: les fondamentalistes l'ont fait refluer. Les ordres de 6\,000 et de 1\,500 actions s'achèvent à
2.6 (1.9) et 2.4 (0.9) ticks, à l'intérieur de leurs marges d'erreur mutuelles. Six séances suffisent à trancher pour le plus grand ordre, mais pas pour la courbe~; la loi en racine carrée a été observée
sur des centaines de milliers de métaordres réels, comme chez Bacry, Iuga, Lasnier et Lehalle (2015), et dans les profils de négociation de Moro et coauteurs (2009).

\begin{omfigure}
\begin{tikzpicture}
\begin{axis}[width=9.5cm,height=5.2cm,xlabel={\small secondes depuis le début du métaordre},ylabel={\small impact (en ticks)},
  tick label style={font=\footnotesize},xmin=-50,xmax=1250,ymin=-6,ymax=12,grid=major,grid style={gray!25},
  xticklabel style={/pgf/number format/.cd,fixed,1000 sep={\,}},
  legend columns=3,legend style={font=\footnotesize,at={(0.5,-0.25)},anchor=north,draw=none,fill=none,
  /tikz/every even column/.append style={column sep=8pt}},table/col sep=comma]
\fill[gray!12] (axis cs:0,-6) rectangle (axis cs:600,12);
\addplot[omDef,thick,mark=*,error bars/.cd,y dir=both,y explicit] table[x=h,y=m1500,y error=e1500]{figdata/microstructure/11-the-empirics-of-price-impact/meta.csv};
\addplot[omThm,thick,mark=square*,error bars/.cd,y dir=both,y explicit] table[x=h,y=m6000,y error=e6000]{figdata/microstructure/11-the-empirics-of-price-impact/meta.csv};
\addplot[black,very thick,mark=diamond*,error bars/.cd,y dir=both,y explicit] table[x=h,y=m24000,y error=e24000]{figdata/microstructure/11-the-empirics-of-price-impact/meta.csv};
\legend{1\,500 actions,6\,000 actions,24\,000 actions}
\end{axis}
\end{tikzpicture}

\omcaption{Impact contrefactuel des métaordres dans le marché simulé~: chaque séance avec le métaordre moins la même séance sans lui, moyenne et erreur type sur six séances~;
la bande grisée correspond à l'exécution. Données~: \texttt{mx\_impact.counterfactual}.}\label{fig:mx:the-empirics-of-price-impact:meta}
\end{omfigure}

\section{La loi en racine carrée}

\begin{definition}[Préfacteur d'impact]\label{def:mx:the-empirics-of-price-impact:prefactor}
Dans la loi en racine carrée $I=Y\sigma\,(Q/V)^{\delta}$ avec $\delta\approx\tfrac12$, où $I$ est l'\omterm{def:mx:the-empirics-of-price-impact:peak}{impact maximal} en fraction du cours, $\sigma$ la
volatilité journalière, $Q$ la taille du métaordre et $V$ le volume journalier, le \emph{préfacteur d'impact}\index{préfacteur d'impact} $Y$ est la constante sans dimension, d'ordre
un, qui fixe le niveau.
\end{definition}

Avec $Y=0.8$ et $\sigma=2~\%$, acheter 1~\% du volume de la journée coûte environ $0.8\times0.02\times\sqrt{0.01}=16$ points de base d'impact, et 4~\% environ 32~: quatre
fois la taille, deux fois l'impact. Toth et coauteurs (2011) ont expliqué la loi par un carnet dont l'offre et la demande moyennes ont une forme en V et s'annulent au
cours, de sorte que les gros métaordres doivent être fractionnés et que la liquidité locale est anormalement faible~; Zarinelli, Treccani, Farmer et Lillo (2015) ont constaté que, sur les
données américaines, la racine carrée s'ajuste sur environ deux ordres de grandeur de taille et un logarithme sur presque cinq. Loeb (1983) avait tabulé des coûts de transaction croissant
avec la taille de l'ordre, bien avant.

Les registres de l'entreprise sont le bon endroit pour l'ajuster~: de nombreux métaordres sur différentes actions, avec leurs tailles, leurs durées, leurs volatilités et leurs exécutions. Le chapitre
les simule~: 2\,000 métaordres sur 50 actions, des tailles allant de 0.01~\% à 5~\% du volume journalier, des durées égales au dixième du volume, une vraie loi $Y=0.8$,
$\delta=\tfrac12$, et un bruit de cours de $\sigma\sqrt T$ sur chacun. \texttt{firm.impactfit.fit\_power} moyenne l'impact sur les tailles à l'intérieur de classes (un impact
isolé peut être négatif, donc les logarithmes des impacts individuels ne servent à rien), ajuste le logarithme des moyennes par classe sur le logarithme de la taille (comme dans le livre~7
\texttt{firm.tcost.fit\_impact}), et applique un bootstrap sur des actions entières pour les intervalles.

\omcode{code/firm/impactfit/firm_impactfit.py}{83}{103}{Exposant et préfacteur de la loi en racine carrée, avec des intervalles à 95~\% obtenus par rééchantillonnage d'actions entières.}\label{lst:mx:the-empirics-of-price-impact:fit}

L'ajustement donne $\delta=0.48$ (intervalle à 95~\% de 0.39 à 0.56) et $Y=0.85$ (de 0.55 à 1.66)~; avec l'exposant fixé à un demi, $Y=0.81$. Les données identifient
l'exposant à environ $\pm0.09$ près et le préfacteur à un facteur deux près, à partir de 2\,000 ordres dont la vraie loi est contenue.

\section{Pièges de la mesure}

\begin{definition}[Contamination par l'alpha]\label{def:mx:the-empirics-of-price-impact:alpha}
La \emph{contamination par l'alpha}\index{contamination par l'alpha} est la part d'un impact mesuré qui correspond au mouvement de cours que le propriétaire de l'ordre attendait de toute façon~: des ordres
envoyés à cause d'une prévision de mouvement, avec la prévision, et dont l'impact paraît plus grand qu'il ne l'est.
\end{definition}

Supposons que 30~\% des métaordres de l'entreprise aient été lancés parce qu'une prévision anticipait une dérive de $0.5\sigma\sqrt T$ dans leur sens sur leur durée. Le même
ajustement donne $\delta=0.50$ (de 0.44 à 0.55) et $Y=1.42$ (de 1.01 à 2.04)~; avec l'exposant fixé, $Y=1.31$, soit 62~\% de trop
(la \cref{fig:mx:the-empirics-of-price-impact:sqrt}). L'exposant n'est pas touché, car l'horizon de la prévision est la durée de l'ordre, qui croît avec sa
taille, de sorte que la dérive varie comme l'impact~; seul le niveau est faux. Retirer la prévision de l'impact mesuré de chaque ordre synchronisé rétablit exactement l'ajustement
propre~; une régression de l'impact sur la prévision aurait trouvé un coefficient de 1.54, et non de un, car les ordres plus gros portent à la fois davantage d'impact et davantage
de prévision.

\omcode{code/microstructure/11-the-empirics-of-price-impact/python/mx_impact.py}{94}{112}{Les métaordres de l'entreprise, simulés~: une loi en racine carrée, du bruit de cours, et une part d'ordres synchronisés par une prévision.}\label{lst:mx:the-empirics-of-price-impact:panel}

\begin{omfigure}
\begin{tikzpicture}
\begin{axis}[width=9cm,height=5.4cm,xmode=log,ymode=log,xlabel={\small $Q/V$},ylabel={\small impact $/\,\sigma$},
  tick label style={font=\footnotesize},xmin=8e-5,xmax=0.06,ymin=0.0008,ymax=0.5,grid=major,grid style={gray!25},
  legend columns=2,legend style={font=\footnotesize,at={(0.5,-0.3)},anchor=north,draw=none,fill=none,
  /tikz/every even column/.append style={column sep=8pt}},table/col sep=comma]
\addplot[omDef,only marks,mark=*] table[x=x,y=clean]{figdata/microstructure/11-the-empirics-of-price-impact/sqrt.csv};
\addplot[omThm,only marks,mark=square*] table[x=x,y=timed]{figdata/microstructure/11-the-empirics-of-price-impact/sqrt.csv};
\addplot[omDef,domain=1e-4:0.05,samples=20] {0.81*sqrt(x)};
\addplot[omThm,dashed,domain=1e-4:0.05,samples=20] {1.31*sqrt(x)};
\legend{prévision retirée,avec ordres synchronisés,$0.81\sqrt{Q/V}$,$1.31\sqrt{Q/V}$}
\end{axis}
\end{tikzpicture}

\omcaption{La loi en racine carrée sur 2\,000 métaordres simulés~: impact moyen par classe divisé par la volatilité journalière, en fonction de la taille divisée par le volume journalier, avec la
prévision retirée et avec 30~\% des ordres synchronisés par une prévision, et les ajustements en racine carrée. La pente est la même~; le niveau ne l'est pas. Données~:
\texttt{mx\_impact.panel\_fits}.}\label{fig:mx:the-empirics-of-price-impact:sqrt}
\end{omfigure}

Les autres pièges sont tout aussi courants. Bruit~: le cours se déplace de $\sigma\sqrt T$ pendant la vie d'un ordre, quoi que fasse cet ordre, de sorte que l'impact des petits ordres est
invisible sans en disposer de nombreux, et le mesurer sur un long horizon fixe fait baisser l'exposant. Sélection~: seuls les ordres exécutés sont enregistrés, et les ordres
annulés parce que le cours s'est éloigné manquent. Chevauchement~: les métaordres simultanés d'une même entreprise, ou du marché, ajoutent leur impact à
celui des autres. Rejeu~: un flux de fond incapable de réagir ne montre aucun impact. Chaque estimation publiée est un choix sur ces points, c'est pourquoi
les estimations de $Y$ diffèrent davantage que celles de $\delta$.

\section{Tutoriel~: la racine carrée dans nos propres données}\label{tut:mx:the-empirics-of-price-impact:tut}

\begin{tutorial}
\textbf{Objectif.} Mesurer l'impact à trois échelles dans un marché réactif, puis ajuster la loi en racine carrée sur les métaordres de l'entreprise et la corriger pour leur synchronisation.
\textbf{État final~:} les \cref{fig:mx:the-empirics-of-price-impact:single,fig:mx:the-empirics-of-price-impact:meta,fig:mx:the-empirics-of-price-impact:sqrt}.
\begin{tutsteps}
\item \textbf{Marché.} \texttt{firm\_agentmkt.session(mx\_impact.MARKET, seconds, seed, agents)}~; \texttt{res.tape()} pour les transactions et les cotations.
\item \textbf{Impact isolé et agrégé.} \texttt{firm\_impactfit.response(trades, times, mids)} et \texttt{aggregate(\dots, 10.0, edges)}.
\item \textbf{Métaordres.} \texttt{mx\_impact.counterfactual()}~: chaque session avec et sans \texttt{Metaorder}~; \texttt{replay\_pitfall()} pour le
flux rejoué.
\item \textbf{Loi.} \texttt{panel(seed, alpha\_share)} et \texttt{fit\_power(impact, sigma, participation, clusters=stock)}~;
\texttt{decontaminate(impact, forecast)}~; dessiner avec \texttt{fig\_impact.py}.
\end{tutsteps}
\textbf{Pour aller plus loin.} Donner aux fournisseurs de liquidité une loi de placement plus plate (forme 0) et voir si la réponse à un ordre isolé devient concave~; ajuster le
panel avec un logarithme au lieu d'une puissance et comparer les résidus.
\end{tutorial}

\section{Construction~: estimation de l'impact}\label{bld:mx:the-empirics-of-price-impact:impactfit}

\begin{build}
\bfield{Rôle} Courbes et lois d'impact tirées des exécutions et des données de marché de l'entreprise, pour les modèles d'exécution des chapitres 14 et 15, les estimations avant transaction du
chapitre~19 et le modèle de coûts du livre~7.
\bfield{Interface} \texttt{orders(trades)}, \texttt{response(trades, times, mids, lags, size\_edges)},
\texttt{aggregate(trades, times, mids, interval, edges)}, \texttt{metaorder\_path(times, mids, start, end, sign, grid)}, \texttt{fit\_power(impact, sigma, participation, bins, boot, clusters)},
\texttt{decontaminate(impact, forecast)}. Avec elle, \texttt{firm.agentmkt} (\texttt{PopulationConfig}, \texttt{Population}, \texttt{Metaorder},
\texttt{session}), dont la description complète est celle du chapitre~27.
\bfield{Règles} Les exécutions à un même instant et dans un même sens forment un seul ordre~; réponses par ordre~; ajustements par classes (\texttt{fit\_impact} du livre~7)~; bootstrap par grappes d'actions~;
prévisions retirées ordre par ordre.
\bfield{Tests d'acceptation} \texttt{code/firm/impactfit/tests/}~: ordres fusionnés et une réponse calculée à la main~; la pente d'un \omterm{def:mx:the-empirics-of-price-impact:aggregate}{impact agrégé} linéaire~; l'exposant
et le préfacteur d'une loi injectée à l'intérieur de leurs intervalles. \texttt{code/firm/agentmkt/tests/}~: une population déterministe, un flux exogène partagé, des signes persistants,
un métaordre qui exécute son plan.
\bfield{Pour aller plus loin} Surfaces d'impact en taille et en durée~; noyaux de décroissance après l'exécution (chapitre~12)~; la loi logarithmique.
\end{build}

\begin{omsources}
\item T.~F.~Loeb, ``Trading cost: the critical link between investment information and results'', \emph{Financial Analysts Journal} 39(3), 1983.
\item E.~Moro, J.~Vicente, L.~G.~Moyano, A.~Gerig, J.~D.~Farmer, G.~Vaglica, F.~Lillo and R.~N.~Mantegna, ``Market impact and trading profile of hidden
orders in stock markets'', \emph{Physical Review E} 80, 2009.
\item B.~Tóth, Y.~Lempérière, C.~Deremble, J.~de~Lataillade, J.~Kockelkoren and J.-P.~Bouchaud, ``Anomalous price impact and the critical nature of
liquidity in financial markets'', \emph{Physical Review X} 1, 2011.
\item E.~Bacry, A.~Iuga, M.~Lasnier and C.-A.~Lehalle, ``Market impacts and the life cycle of investors orders'', \emph{Market Microstructure and
Liquidity} 1(2), 2015.
\item E.~Zarinelli, M.~Treccani, J.~D.~Farmer and F.~Lillo, ``Beyond the square root: evidence for logarithmic dependence of market impact on size and
participation rate'', \emph{Market Microstructure and Liquidity} 1(2), 2015.
\end{omsources}

\section{Exercices}

\begin{exercise}[$\star$]\label{exo:mx:the-empirics-of-price-impact:1}
Avec $Y=0.8$, $\delta=\tfrac12$ et une volatilité journalière de 2~\%, quel impact la loi en racine carrée donne-t-elle pour l'achat de 1~\% et de 4~\% du volume de la journée~?
\end{exercise}

\begin{exercise}[$\star$]\label{exo:mx:the-empirics-of-price-impact:2}
Dans le marché simulé, le milieu se déplace de 2.1 ticks pour mille actions d'achat net sur dix secondes. De combien doivent 600 actions d'achat net le déplacer, et
pourquoi la relation n'est-elle qu'approximativement linéaire~?
\end{exercise}

\begin{exercise}[$\star$]\label{exo:mx:the-empirics-of-price-impact:3}
Le métaordre de 24\,000 actions a fait bouger le cours de 7.8 ticks à sa fin, puis de $-1.5$ cinq minutes plus tard. Quelle part de son impact a été résorbée, et qu'est-ce qui
l'a causée sur le marché~?
\end{exercise}

\begin{exercise}[$\star\star$]\label{exo:mx:the-empirics-of-price-impact:4}
Pourquoi un prévisionnel dont l'horizon est la durée du métaordre biaise-t-il le préfacteur mais pas l'exposant~?
\end{exercise}

\begin{exercise}[$\star\star$]\label{exo:mx:the-empirics-of-price-impact:5}
Pourquoi les intervalles calculés par rééchantillonnage portent-ils sur des actions entières plutôt que sur des ordres isolés~?
\end{exercise}

\begin{exercise}[$\star\star$]\label{exo:mx:the-empirics-of-price-impact:6}
Pourquoi un fond de marché rejoué ne montre-t-il aucun impact, et qu'est-ce que cela implique pour les backtests sur données rejouées~?
\end{exercise}

\begin{exercise}[$\star\star\star$]\label{exo:mx:the-empirics-of-price-impact:7}
\emph{Programmation.} Relancer le panel en fixant la durée de chaque ordre à un dixième de jour, avec et sans les ordres datés. Ajuster la loi et expliquer ce que
vous trouvez.
\end{exercise}

\begin{exercise}[$\star\star\star$]\label{exo:mx:the-empirics-of-price-impact:8}
\emph{Trouver la faille.} «~Nos exécutions donnent un exposant de 0.5, exactement la valeur publiée, donc notre modèle d'impact est juste et notre alpha n'y figure pas.~»
\end{exercise}

\section{Problème~: la racine carrée dans nos propres données}

\begin{problem}[{Problème du week-end --- la racine carrée dans nos propres données}]\label{pb:mx:the-empirics-of-price-impact:1}
Le desk d'exécution souhaite un modèle d'impact pré-trade et demande à la recherche de l'ajuster sur les métaordres propres de la société. Mesurer, ajuster et vérifier.

\medskip\noindent\textbf{Partie I --- Petites échelles.}
\begin{enumerate}
\item Décrire le marché simulé et ce qui y réagit.
\item Donner la réponse à un \omterm{def:mx:the-limit-order-book:orders}{ordre au marché} selon sa taille, après un ordre puis après cinquante.
\item Pourquoi la réponse des ordres d'un lot grandit-elle avec le décalage~?
\item Donner l'\omterm{def:mx:the-empirics-of-price-impact:aggregate}{impact agrégé} par millier d'actions sur dix secondes.
\end{enumerate}

\medskip\noindent\textbf{Partie II --- Métaordres.}
\begin{enumerate}[resume]
\item Comment mesure-t-on un impact contrefactuel dans un simulateur, et que retire-t-il~?
\item Que montre le fond rejoué, et pourquoi~?
\item Donner l'impact de l'ordre de 24\,000 actions à mi-parcours, à la fin et cinq minutes après.
\item Pourquoi six séances ne suffisent-elles pas à résoudre la \omterm{def:mx:the-empirics-of-price-impact:curve}{courbe d'impact}~?
\end{enumerate}

\medskip\noindent\textbf{Partie III --- La loi.}
\begin{enumerate}[resume]
\item Énoncer la loi en racine carrée et calculer l'exercice~1.
\item Comment la loi est-elle ajustée sur le panel, et pourquoi par classes~?
\item Donner l'exposant et le préfacteur avec leurs intervalles.
\item Que trouvent Toth et ses coauteurs, et Zarinelli et ses coauteurs~?
\end{enumerate}

\medskip\noindent\textbf{Partie IV --- Pièges et verdict.}
\begin{enumerate}[resume]
\item Définir la \omterm{def:mx:the-empirics-of-price-impact:alpha}{contamination par l'alpha}.
\item Donner l'ajustement avec 30~\% des ordres datés par un prévisionnel.
\item Pourquoi l'exposant est-il sans biais ici~?
\item Comment retire-t-on le prévisionnel, et que renvoie une régression sur celui-ci~?
\item Que se passe-t-il avec un horizon long fixe (exercice~7)~?
\item Énumérer trois autres pièges de mesure.
\item Énoncer le \emph{résultat nommé}~: l'exposant et le préfacteur de la loi en racine carrée ajustée sur les métaordres simulés de la société, avec leurs intervalles de confiance,
et leur biais lorsque le calendrier des ordres porte un prévisionnel.
\item En une phrase~: que doit utiliser le modèle pré-trade du desk~?
\end{enumerate}
\end{problem}

\section{Questions d'entretien}

\begin{interviewq}[$\star$ \iqroles{trader, chercheur}]\label{iq:mx:the-empirics-of-price-impact:1}
Énoncez la loi d'impact en racine carrée et ce que signifie chacun de ses termes.
\end{interviewq}

\begin{interviewq}[$\star\star$ \iqroles{chercheur}]\label{iq:mx:the-empirics-of-price-impact:2}
Comment estimeriez-vous l'impact des métaordres de votre société, et qu'est-ce qui biaiserait l'estimation~?
\end{interviewq}

\begin{interviewq}[$\star\star$ \iqroles{chercheur}]\label{iq:mx:the-empirics-of-price-impact:3}
Pourquoi l'impact est-il concave en fonction de la taille~? Donnez une explication.
\end{interviewq}

\begin{interviewq}[$\star\star$ \iqroles{trader}]\label{iq:mx:the-empirics-of-price-impact:4}
Votre backtest sur données de marché rejouées n'affiche aucun coût pour un ordre important au-delà de la fourchette. Qu'est-ce qui cloche~?
\end{interviewq}

\begin{interviewq}[$\star\star$ \iqroles{chercheur}]\label{iq:mx:the-empirics-of-price-impact:5}
Quelle part de l'impact d'un métaordre est permanente, et de quoi dépend la réponse~?
\end{interviewq}

\begin{interviewq}[$\star\star\star$ \iqroles{chercheur}]\label{iq:mx:the-empirics-of-price-impact:6}
Les ordres de votre équipe alpha montrent deux fois l'impact de ceux de l'équipe indicielle, pour une même taille. Leur exécution est-elle moins bonne~?
\end{interviewq}
